\chapter*{Przedmowa}

\phantomsection

\lettrine[lines=5, lhang=0.1]{G}{dy} byłem dzieckiem, lubiłem słuchać
opowiadań moich dziadków, o~ich życiu, pracy i~rodzinie, nie były to historie
o~wielkich bitwach, podróżach czy~odkryciach, bo woleli oni przekazywać raczej
swoje niż cudze przeżycia. Dziadek Stasiek i~babcia Jadzia
najbardziej lubili snuć swoje drobne historie w~trakcie wspólnej pracy
w~ogródku, na~polu, na~podwórku, czy~w~ich~rodzinnym domu, który razem, 
z~mozołem wznieśli pracą swoich skromnych rąk. Dom ten, istniał niegdyś pod
adresem: \emph{Desno 14}, był to mój pierwszy dom rodzinny, ale~teraz,
bez~nich, zostały z~niego już tylko puste mury.

Wspólny posiłek i~wspólna praca były od zawsze kluczem do~opowiadania historii
w~mojej rodzinie, gdyż dziadkowie Gomulscy należeli do~tego przedwojennego
pokolenia, które nigdy nie marnowało sił bez potrzeby. Gdy tylko mieli choć
odrobinę więcej energii życiowej to~zawsze, pomimo licznych przypadłości
i~schorzeń, znajdowali coś do~naprawienia, przeniesienia, ugotowania,
posprzątania, czy wypielenia. A~ja, wraz z~moim starszym bratem Krzyśkiem oraz
bratem ciotecznym Łukaszem, z~racji, iż mieszkaliśmy blisko nich, bardzo
często im w~tych pracach mimowolnie towarzyszyliśmy. Widzę to teraz w~moich
rodzicach Andrzeju i~Bożenie, że ten przedwojenny etos ciężkiej pracy dziadków,
dalej jest w~nich bardzo żywy.

Skłamałbym gdybym napisał, że moi dziadkowie nie mieli nigdy żadnego 
\emph{hobby} (dziadek podobno całkiem nieźle grywał niegdyś na harmonijce
ustnej), ale~przynajmniej za mojego świadomego życia, każdy swój gram
nadmiarowej energii poświęcali na pracę nad swoim rodzinnym poletkiem, a~gdy
nie mieli siły, to już tylko przesiadywali -- przy~oknie, przy~telewizorze,
przy~radiu, czy~na ławce w~ogródku, poczciwie doglądając otaczający ich~świat.
Wyjątek od tej reguły stanowiło jedynie grzybobranie, to bez wątpienia była
pasja babci Jadzi, praktycznie do~końca jej życia, czy~miała siłę, czy~też~nie,
a~ja i~tu byłem jej wiernym towarzyszem i~słuchaczem\footnote{Tradycję
wędrowania za grzybami po deśnińskich lasach, kępach i~brzezinach kultywuję,
z~sukcesami, do~dnia dzisiejszego, niestety już nie w~towarzystwie babci.}.

Prawdę mówiąc, większości historii opowiadanych przez moich dziadków wtedy nie
rozumiałem, wiele pewnie zapomniałem, ale~cześć z~nich wróciła do~mnie już
w~trakcie pisania tej książki -- by nabrać wreszcie pełnego smaku. Ta książka
nie będzie jednak biografią moich dziadków, tylko próbą zrozumienia tego co
ich ukształtowało, jak do~tego doszło, że oboje urodzili się w~małej
podwarszawskiej wsi, jako~dzieci z~ubogich chłopskich rodzin. Jako~dobrą
ilustrację losów ich~przodków, zdecydowałem się wybrać historię rodziny
dziadka Staśka, gdyż była ona dużo bardziej związana z~naszym rodzinnym Desnem,
można o~niej odnaleźć więcej informacji w~dokumentach źródłowych oraz była
również istotną częścią historii matki mojej babci -- Magdaleny Poziemskiej
(z~domu Zawadzkiej, z~pierwszego małżeństwa Gomulskiej). Rodziny Poziemskich
i~Zawadzkich na Deśnie wygasły niestety wraz ze śmiercią mojej prababki,
w~lutym 1970~roku, oraz babci, w~marcu 2014~roku\footnote{Babcia Jadzia zawsze
powtarzała, że marzec i~listopad to najtrudniejsze do~przeżycia miesiące dla
ludzi starszych i~oboje z~dziadkiem zmarli niestety właśnie w~marcu --
dziadek 9~lat po babci w~2023~roku, dożywając aż 96~lat i~333~dni.}, to też
nigdy nie poznałem żadnych krewnych ze strony tych rodzin, a~krewnych
ze~strony rodziny Gomulskich mam całe mnóstwo. Być może kiedyś zdecyduję się
napisać też coś więcej o~historii rodzin Poziemskich oraz~Zawadzkich,
ale~w~tej książce ograniczę się jedynie do przedstawienia pokrótce
ich~genealogii.

Większość historii, które poznałem w~trakcie tej genealogicznej przygody, było
dla mnie zupełnie nowych, gdyż dziadkowie prawdopodobnie sami ich~nie znali,
albo~postanowili się nimi ze mną nie dzielić, albo~przepadły one w~czeluściach
mojej pamięci oraz pamięci moich najbliższych. Ta książka powstała by zebrać
i~usystematyzować te wszystkie historie: te opowiadane przez moich dziadków
oraz te zapisane tuszem w~starych dokumentach. Powstała również po to by
odpowiedzieć na liczne pytania, które nurtowały mnie, czasem już od wczesnych
lat mojego dzieciństwa:

\emph{Skąd się wzięło nasze nazwisko?} \\
\emph{Skąd pochodzą moi przodkowie i~jak wyglądała ich~przeszłość?} \\
\emph{Czy~dziadek Stasiek był najdłużej żyjącym członkiem rodziny
Gomulskich?} \\
\emph{Czy~rzeczywiście większość moich sąsiadów jest ze mną tak blisko
spokrewniona?} \\
\emph{Skąd się wzięła nazwa mojej rodzinnej miejscowości? Kiedy została
założona i~przez kogo?}

Na większość tych kwestii nie mogłem uzyskać satysfakcjonujących mnie
wyjaśnień, ani~od~dziadków, ani~od~innych krewnych, ani~od nauczycieli, którzy
albo nie znali na nie odpowiedzi, albo~mieli różne, często sprzeczne
wytłumaczenia. Ta~ciekawość doprowadziła mnie w~czerwcu 2022~roku
do~rozpoczecia tej~fascynującej kilkuletniej podróży genealogicznej -- w~głąb
historii mojej rodziny. Podróży, która pozwoliła mi odkryć wiele niesamowitych
opowieści o~moich przodkach, ich życiu, pracy, a~także o~tym, jak bardzo
zmieniła się Polska na przestrzeni ostatnich trzech wieków. Podróży, która
pozwoliła mi zrozumieć, że historia mojej rodziny jest w~gruncie rzeczy
zapisem losów wielu innych zwykłych rodzin w~Polsce, które również
doświadczyły podobnych wyzwań, zmian, sukcesów oraz porażek.

Teraz, gdy moi dziadkowie już nie żyją, pryszedł czas, żebym zebrał te
wszystkie opowieści i~sam je przekazał, by nie zaginęły i~nie zostały w~pełni
zapomniane (tak, Haniu i~Gusiu, jeśli to czytacie, to często myślałem o Was,
gdy pisałem tę książkę). Zwyczaj wspólnej pracy całych rodzin już niestety
powoli zanika, więc ten pas transmisyjny przenoszący rodzinne historie
i~tradycje, z~pokolenia na pokolenie, powoli koroduje. Dlatego właśnie, poza
wcześniej wspomnianymi powodami, zdecydowałem się na napisanie tej książki,
ale~również dlatego, że zawsze łatwiej mi było odnaleźć odpowiednie słowa
pisząc niż mówiąc, za~co~nieraz, do~dziś, podpadam mojej ukochanej żonie,
Ewie.

Piszę więc ten wstęp z~krawędzi szpitalnego łóżka, obserwując krople powoli
spływające z~kolejnej dawki mojej pierwszej chemioterapii, wciąż żywiąc
naiwną nadzieję, że nie okaże się on, z~biegiem czasu, zbyt ckliwy
i~melancholijny. Zapewniam jednak, iż dalsze rozdziały tej książki, skupiają
się już -- stety lub niestety -- prawie wyłącznie na datach, tabelkach
i~dokumentach źródłowych, bo jest to główny język genealogii. Język
nieprzystępny, dla wielu osób czasami nużący, dlatego też starałem się go
uatrakcyjnić, jak tylko mogłem, aby okazał się bardziej strawny dla kolejnych,
bardziej już cyfrowych pokoleń.

Dedykuję tę książkę moim zmarłym dziadkom, którzy byli dla mnie nie tylko
wspaniałymi opiekunami i~wychowawcami, wspierając w~tym trudnym zadaniu moich
rodziców, ale~byli, przede wszystkim, niesamowitymi nauczycielami, którzy
przekazali mi, często pewnie nieświadomie, wiele cennych lekcji o~życiu,
miłości, pracy, i~relacjach międzyludzkich, które powoli do~mnie przenikały
wraz z~upływem kolejnych lat mojego życia.

Ciebie zaś, drogi Czytelniku, serdecznie zapraszam do~wspólnej podróży w~głąb
historii mojej skromnej rodziny -- Gomulskich z~Desna. \\

\begin{flushright}
\emph{Mateusz Gomulski, 10~kwietnia 2026~roku} \\
\emph{Warszawa ul. Szamocka 6, Szpital Onkologiczny} \\
\emph{Początek Pierwszego Turnusu Leczenia Onkologicznego}
\end{flushright}